% Einheit 2 von 4 – Terme zusammenfassen (bank/terme/e2.jsonl, zusammenbau v0.9)
\einheitenkopf[e2]{Einheit 2 von 4 · Terme zusammenfassen}
\zweigzeile{Hier lernst du, Terme zusammenzufassen · ab Kl. 7, je nach Buch bis Kl. 8 · P10 ×1 · baut auf: Addieren und Subtrahieren negativer Zahlen, Multiplizieren mit Vorzeichen, Dezimalzahlen und einfache Brüche als Vorzahlen, Punkt vor Strich mit Zahlen}
\setcounter{aufgabe}{14}


% erkennung: terme-e2-k1-s0-v1, terme-e2-k1-s0-v2, terme-e2-k1-s0-v3, terme-e2-k1-s0-v4
\begin{kbaufgabe}{Ich erkenne gleichartige Glieder.}
\begin{kbblock}
\kbanweisung{Unterstreiche, was zusammengehört.}
\lbpaar{\kbteil{a)}{$4y$, $7$, $9y$}{}
}{\kbteil{b)}{$6$, $2a$, $8$, $a$}{}}
\end{kbblock}
\begin{kbblock}
\lbpaar{\kbteil{c)}{$5x^2$, $4x$, $3x^2$}{}
}{\kbteil{d)}{$7b$, $2c$, $5b$, $4$}{}}
\end{kbblock}
\end{kbaufgabe}


% erkennung: terme-e2-k2-s0-v1, terme-e2-k2-s0-v2, terme-e2-k2-s0-v3, terme-e2-k2-s0-v4
\begin{kbaufgabe}{Ich kann die Vorzahl eines Gliedes ablesen.}
\begin{kbblock}
\kbanweisung{Welche Vorzahl hat …?}
\lbpaar{\kbteil{a)}{$a$ \quad \lbfeld}{}
}{\kbteil{b)}{$-b$ \quad \lbfeld}{}}
\end{kbblock}
\begin{kbblock}
\lbpaar{\kbteil{c)}{$1{,}5z$ \quad \lbfeld}{}
}{\kbteil{d)}{$-4w$ \quad \lbfeld}{}}
\end{kbblock}
\end{kbaufgabe}


% erkennung: terme-e2-k3-s0-v1, terme-e2-k3-s0-v2, terme-e2-k3-s0-v3, terme-e2-k3-s0-v4
\begin{kbaufgabe}{Ich kann die Glieder eines Terms mit ihrem Vorzeichen aufschreiben.}
\begin{kbblock}
\kbanweisung{Schreibe die Glieder des Terms … mit ihrem Vorzeichen auf.}
\lbpaar{\kbteil{a)}{$5a - 2b + 7$ \quad \lbfeld}{}
}{\kbteil{b)}{$-3x + 6 - y$ \quad \lbfeld}{}}
\end{kbblock}
\begin{kbblock}
\lbpaar{\kbteil{c)}{$2m - 9 - 5k$ \quad \lbfeld}{}
}{\kbteil{d)}{$-6 + 3c - 4d$ \quad \lbfeld}{}}
\end{kbblock}
\end{kbaufgabe}


% leiter: terme-e2-k4-s1-v1, terme-e2-k4-s1-v2, terme-e2-k4-s1-v3, terme-e2-k4-s1-v4, terme-e2-k4-s2-v1, terme-e2-k4-s3-v1, terme-e2-k4-s4-v1
\begin{kbaufgabe}{Ich kann Terme zusammenfassen.}
\begin{kbblock}
\kbanweisung{Fasse zusammen.}
\lbpaar{\kbteil{a)}{$7x + 2x =$ \lbfeld}{}
}{\kbteil{b)}{$7x + 4x =$ \lbfeld}{}}
\end{kbblock}
\begin{kbblock}
\lbpaar{\kbteil{c)}{$7x + 5x =$ \lbfeld}{}
}{\kbteil{d)}{$7x + 8x =$ \lbfeld}{}}
\end{kbblock}
\begin{kbblock}
\lbpaar{\kbteil{e)}{$9x - 4x =$ \lbfeld}{}
}{\kbteil{f)}{$2x + 4x + 3x =$ \lbfeld}{}}
\end{kbblock}
\begin{kbblock}
\lbpaar{\kbteil{g)}{$5x + 3 + 2x =$ \lbfeld}{}
}{\item[]}
\end{kbblock}
\end{kbaufgabe}


% leiter: terme-e2-k4-s5-v1, terme-e2-k4-s6-v1, terme-e2-k4-s7-v1, terme-e2-k4-s8-v1, terme-e2-k4-s9-v1, terme-e2-k4-s10-v1
\begin{kbaufgabe}{Ich kann Terme zusammenfassen. – weiter}
\begin{kbblock}
\kbanweisung{Fasse zusammen.}
\lbpaar{\kbteil{a)}{$4a + 3b + 2a =$ \lbfeld}{}
}{\kbteil{b)}{$y + 4y =$ \lbfeld}{}}
\end{kbblock}
\begin{kbblock}
\lbpaar{\kbteil{c)}{$2x - 7x =$ \lbfeld}{}
}{\kbteil{d)}{$-3x + 8x =$ \lbfeld}{}}
\end{kbblock}
\begin{kbblock}
\lbpaar{\kbteil{e)}{$4x^2 + 5x + 2x^2 =$ \lbfeld}{}
}{\kbteil{f)}{$1{,}5x + 3{,}5x =$ \lbfeld}{}}
\end{kbblock}
\end{kbaufgabe}


% pruefung: terme-e2-k4-s11-v1
\begin{kbaufgabe}{Ich kann auch Aufgaben aus der Prüfung zum Zusammenfassen lösen.}
\begin{kbblock}
\kbteil{}{Vereinfache den Term $8a - 2a^2 - 5a$ und berechne seinen Wert für $a = 3$.}{P10 ’25 G}
\item[]\kbraum{3}
\end{kbblock}
\end{kbaufgabe}


% leiter: terme-e2-k5-s1-v1, terme-e2-k5-s1-v2, terme-e2-k5-s1-v3, terme-e2-k5-s1-v4, terme-e2-k5-s2-v1, terme-e2-k5-s3-v1, terme-e2-k5-s4-v1, terme-e2-k5-s5-v1, terme-e2-k5-s6-v1
\begin{kbaufgabe}{Ich kann Terme malnehmen.}
\begin{kbblock}
\kbanweisung{Multipliziere.}
\lbpaar{\kbteil{a)}{$5 \cdot 2x =$ \lbfeld}{}
}{\kbteil{b)}{$5 \cdot 3x =$ \lbfeld}{}}
\end{kbblock}
\begin{kbblock}
\lbpaar{\kbteil{c)}{$5 \cdot 6x =$ \lbfeld}{}
}{\kbteil{d)}{$5 \cdot 7x =$ \lbfeld}{}}
\end{kbblock}
\begin{kbblock}
\lbpaar{\kbteil{e)}{$4x \cdot 5 =$ \lbfeld}{}
}{\kbteil{f)}{$4y \cdot 2y =$ \lbfeld}{}}
\end{kbblock}
\begin{kbblock}
\lbpaar{\kbteil{g)}{$(-4) \cdot 5x =$ \lbfeld}{}
}{\kbteil{h)}{$4a \cdot 5b =$ \lbfeld}{}}
\end{kbblock}
\begin{kbblock}
\lbpaar{\kbteil{i)}{$2 \cdot 5x \cdot 3 =$ \lbfeld}{}
}{\item[]}
\end{kbblock}
\end{kbaufgabe}


% pruefung: terme-e2-k5-s7-v1
\begin{kbaufgabe}{Ich kann auch Produkte mit Vorzeichen und zwei Variablen berechnen.}
\begin{kbblock}
\kbanweisung{Multipliziere.}
\kbteil{}{$(-6a) \cdot 2b$}{}
\kbantwort{= \kbfeld}
\end{kbblock}
\end{kbaufgabe}


% pflicht: terme-e2-k6-s1-v1, terme-e2-k6-s1-v2, terme-e2-k6-s1-v3
\begin{kbaufgabe}{Ich finde den Fehler beim Zusammenfassen.}
\begin{kbblock}
\kbteil{a)}{Mia soll den Term $5a + 3 + 2a$ zusammenfassen. Sie rechnet so: \rechnung{5a + 3 + 2a &= 10a} Finde den Fehler und rechne richtig.}{}
\item[]\kbraum{2}
\end{kbblock}
\begin{kbblock}
\kbteil{b)}{Tom soll den Term $b + 8b - 3$ zusammenfassen. Er rechnet so: \rechnung{b + 8b - 3 &= 9b - 3} Prüfe, ob Tom richtig gerechnet hat.}{}
\item[]\kbraum{2}
\end{kbblock}
\begin{kbblock}
\kbteil{c)}{Jana hat vier Terme zusammengefasst. Welche Ergebnisse können nicht stimmen? Begründe, ohne genau zu rechnen.}{}
\kbfrage{$4x + 5x = 9x^2$ \\ $6y - 11y = 5y$ \\ $2a + 7a = 9a$ \\ $3b \cdot 2b = 6b$}
\item[]\kbraum{2}
\end{kbblock}
\end{kbaufgabe}


% pflicht: terme-e2-k6-s2-v1, terme-e2-k6-s2-v2, terme-e2-k6-s2-v3
\begin{kbaufgabe}{Ich kann begründen, ob man Glieder zusammenfassen kann.}
\begin{kbblock}
\kbteil{a)}{Begründe, ob man den Term $4x + 4y$ zusammenfassen kann.}{}
\item[]\kbraum{3}
\end{kbblock}
\begin{kbblock}
\kbteil{b)}{Entscheide bei jeder Aussage, ob sie wahr oder falsch ist. Begründe, ohne genau zu rechnen.}{}
\kbfrage{(1) Zwei Glieder mit derselben Variablen und derselben Hochzahl kann man immer zusammenfassen. \\ (2) Zwei Glieder mit derselben Vorzahl kann man immer zusammenfassen. \\ (3) Beim Zusammenfassen entsteht nie eine neue Hochzahl.}
\item[]\kbraum{3}
\end{kbblock}
\begin{kbblock}
\kbteil{c)}{Ole fasst $3y - 8y$ zusammen. Er sagt: „Das ergibt $5y$, denn 8 minus 3 ist 5.“ Begründe, ob Ole recht hat.}{}
\item[]\kbraum{3}
\end{kbblock}
\end{kbaufgabe}


% pflicht: terme-e2-k6-s3-v1, terme-e2-k6-s3-v2, terme-e2-k6-s3-v3
\begin{kbaufgabe}{Ich kann Terme zu Figuren aufstellen und zusammenfassen.}
\begin{kbblock}
\kbteil{a)}{Die Figur zeigt ein Rechteck. Die Maße sind in cm angegeben. Schreibe einen Term für den Umfang.}{}
\kbdarunter{\viereck[seiten={3x,2,3x,2}]{(0,0)}{(6,0)}{(6,2)}{(0,2)}}
\kbantwort{\kbfeld}
\end{kbblock}
\begin{kbblock}
\kbteil{b)}{Die Figur zeigt ein Dreieck. Die Maße sind in cm angegeben. Schreibe einen Term für den Umfang.}{}
\kbdarunter{\dreieck{(0,0)}{(4.8,0)}{(1.65,1.74)}{3a}{2a}{4a}{}{}{}}
\kbantwort{\kbfeld}
\end{kbblock}
\begin{kbblock}
\kbteil{c)}{Der Term $3a + 3a + 3a$ gibt den Umfang einer Figur in cm an. Beschreibe eine passende Figur in Worten, ohne sie zu zeichnen. Fasse den Term dann zusammen.}{}
\item[]\kbraum{2}
\end{kbblock}
\end{kbaufgabe}


% pflicht: terme-e2-k6-s4-v1, terme-e2-k6-s4-v2, terme-e2-k6-s4-v3
\begin{kbaufgabe}{Ich kann Terme im Alltag aufstellen und zusammenfassen.}
\begin{kbblock}
\kbteil{a)}{Ein Kiosk verkauft Brötchen zu je $p$ Euro. Am Montag verkauft er 45 Stück, am Dienstag 38 und am Mittwoch 52. Stelle einen Term für die Einnahmen der drei Tage auf und fasse ihn zusammen. Der Kiosk will in den drei Tagen mindestens 50 € mit Brötchen einnehmen. Schafft er das bei $p = 0{,}40$?}{}
\item[]\kbraum{3}
\end{kbblock}
\begin{kbblock}
\kbteil{b)}{Ein Beet ist $s$ Meter breit und doppelt so lang. Um das Beet kommt ein Zaun. Stelle einen Term für die Länge des Zauns auf und fasse ihn zusammen. Wie viele Meter Zaun braucht man bei $s = 8$?}{}
\item[]\kbraum{2}
\end{kbblock}
\begin{kbblock}
\kbteil{c)}{$n$ Schüler fahren auf Klassenfahrt. Jeder zahlt 85 € für die Unterkunft, 24 € für Eintritte und 36 € für die Bahn. Stelle einen Term für den Gesamtbetrag auf und fasse ihn zusammen. Wie viel Euro kommen bei 26 Schülern zusammen?}{}
\item[]\kbraum{2}
\end{kbblock}
\end{kbaufgabe}
